\section[Appendix - Discovery Workbook Instrument]{Discovery Workbook Instrument}\label{app:workbook}

This appendix reproduces the complete Round~1 discovery workbook distributed to
participants in January 2026. The workbook comprises seven parts (A--G) that
progress from context setting through practice discovery to reflective
sensemaking.

\textbf{Part A: Context Setting}

This opening section establishes the participant's GenAI context and experience
level:
\begin{itemize}
    \item Which GenAI tools do you currently have access to? (e.g., ChatGPT, Claude,
          Copilot, Gemini, internal tools)
    \item How long have you been experimenting with GenAI in your EA work?
    \item What triggered you to start using GenAI for Reference Architecture development?
    \item Describe your organizational context for EA practice
\end{itemize}

\textbf{Part B: Practice Archaeology}

The core discovery section employs multiple prompts to excavate practices from
different angles:

\textbf{B.1 Genesis Moments}\\
``Think back to the last 3 Reference Architectures you've worked on. Walk through a typical RA development effort.
At what points did you think `maybe GenAI could help here'? What did you actually try?''

\textbf{B.2 Practice Inventory}\\
``Now systematically review your GenAI use. For EACH distinct way you use GenAI, please document:''
\begin{enumerate}[label=B.2.\arabic*]
    \item What do you call this practice? (give it a name)
    \item What problem does it solve?
    \item Describe exactly what you do (step-by-step process)
    \item What specific prompts or techniques have you developed?
    \item How did you discover or develop this approach?
    \item What percentage of the time does it work effectively?
    \item What are the typical failure modes?
\end{enumerate}

Participants documented two to five core practices using this structure.

\textbf{Part C: Lifecycle Positioning}

Participants receive a standard RA lifecycle model \textit{without any GenAI
    annotations}:
\begin{enumerate}
    \item Stakeholder Analysis \& Requirements Gathering
    \item Current State Analysis
    \item Pattern Identification \& Abstraction
    \item Target State Design
    \item Component Specification
    \item Documentation \& Communication
    \item Validation \& Review
    \item Evolution \& Maintenance
\end{enumerate}

For each phase, participants answer:
\begin{itemize}
    \item Which of your named practices apply in this phase?
    \item What percentage of this phase involves GenAI assistance?
    \item What tasks MUST remain human-performed? Why?
    \item What surprised you about GenAI's capabilities or limitations here?
\end{itemize}

\textbf{Part D: Critical Incidents}

``Share 3-5 specific stories where GenAI significantly impacted your RA work (positive or negative).''

For each incident, participants provide:
\begin{itemize}
    \item \textbf{Context}: What were you trying to accomplish?
    \item \textbf{Action}: Exactly what did you do with GenAI?
    \item \textbf{Result}: What happened? (intended and unintended outcomes)
    \item \textbf{Insight}: What did this teach you about using GenAI for RA work?
    \item \textbf{Evidence}: Can you share specific prompts and outputs? (anonymized as needed)
\end{itemize}

\textbf{Part E: Impact Reflection}

\textbf{E.1 Time and Effort Analysis}
\begin{enumerate}[label=E.1.\arabic*]
    \item Estimate hours spent on your last RA with vs. without GenAI
    \item Count iteration cycles needed with GenAI assistance
    \item Identify where GenAI saved time
    \item Identify where GenAI required additional time investment
\end{enumerate}

\textbf{E.2 Quality and Completeness Assessment}\\
``Comparing GenAI-assisted RAs to traditionally developed ones:''
\begin{enumerate}[label=E.2.\arabic*]
    \item What aspects show improved quality? Why?
    \item What aspects show degraded quality? Why?
    \item What is simply different rather than better or worse?
\end{enumerate}

\textbf{E.3 Role Evolution}\\
``How has YOUR role as an EA practitioner changed?''
\begin{enumerate}[label=E.3.\arabic*]
    \item Skills you use more frequently now
    \item Skills you use less frequently now
    \item New skills you've had to develop
    \item What would happen to your practice if GenAI disappeared tomorrow?
\end{enumerate}

\textbf{Part F: Challenges and Workarounds}

``Document obstacles you've encountered using GenAI for Reference Architecture development.''

For each challenge:
\begin{itemize}
    \item Describe the specific problem
    \item Explain your workaround (if any)
    \item Identify what capability you wish GenAI possessed
    \item Note any red flags you've learned to recognize
\end{itemize}

\textbf{Part G: Emergent Wisdom}

``Based on your experience so far:''
\begin{itemize}
    \item What essential advice would you give an EA colleague starting with GenAI
          tomorrow?
    \item What critical mistakes should they avoid?
    \item What's your most counterintuitive discovery about GenAI in RA work?
    \item Where do you predict GenAI use in RA development will be in 12 months?
\end{itemize}
